\chapter{The Big Picture}
\label{ch:overview}

Before we describe a single circuit, we need a mental model of the chip that
will contain it. This chapter answers three questions: what is inside an FPGA,
what is on our board, and which tools bring a description in SystemVerilog onto
the chip.

\section{What an FPGA Is}
\label{sec:fpga}

A processor, say the CPU of your laptop, is a fixed circuit. It was designed
once, and all chips of this type contain the same gates connected by the same
wires. What changes is the program it runs.

An FPGA, a \term{field programmable gate array}, is different. It is a chip
that consists of a large number of small, identical building blocks and of
wires between them, and both are \emph{configurable}: which Boolean function a
block computes, and which blocks are connected to which. The configuration is
loaded into the chip when it is switched on. So the circuit itself is what we
change, not a program running on it. If we configure an FPGA as a processor,
then this processor can in turn run programs, and that is exactly what we do
in the last step.

The building blocks of our FPGA, a Lattice iCE40UP5K, come in a few kinds:
\begin{description}
\item[Logic cells.] The workhorse. A logic cell consists of a lookup table with
  four inputs, a flip-flop, and a little extra logic for adding numbers. The
  iCE40UP5K has 5280 of them.
\item[Routing.] Wires between the logic cells, with switches that decide which
  wires are connected. Most of the area of an FPGA is routing, not logic.
\item[I/O pins.] The connections to the outside world, here buttons, LEDs, the
  display and the USB chip on the board.
\item[Memory blocks.] Ready-made memory, because building memory from
  flip-flops would use up the logic cells quickly. We meet them in
  Chapter~\ref{ch:step07}.
\end{description}
The lookup table is the part that needs explaining, and it is simpler than its
name suggests.

\section{A Lookup Table Is a Multiplexer}
\label{sec:lut}

You know the multiplexer from CircuitVerse: it has data inputs, control inputs,
and one output, and the control bits select which data input appears at the
output. A multiplexer with four control bits selects one of $2^4 = 16$ data
inputs.

Now take such a 16-to-1 multiplexer and do not connect its data inputs to
signals. Instead, fix each of them to a constant 0 or 1 when the chip is
configured. These 16 constants are the \term{configuration bits} of the
lookup table, and the four control inputs are its actual inputs
(Figure~\ref{fig:lut}). This is a \term{lookup table}, or \term{LUT} for short.

\begin{figure}[ht]
\centering
\begin{tikzpicture}[x=1mm,y=1mm]
  % the multiplexer
  \draw[thick,fill=green!8] (0,0) -- (0,64) -- (18,54) -- (18,10) -- cycle;
  \node[font=\small] at (9,32) {MUX};
  \foreach \k in {0,...,15} {
    \pgfmathsetmacro{\y}{60-3.7*\k}
    \draw[wire] (-12,\y) -- (0,\y);
    \node[font=\scriptsize,anchor=east] at (-12,\y) {$t_{\k}$};
  }
  \node[note,anchor=east] at (-20,32) {16 configuration\\bits, set when\\the FPGA is\\configured};
  \draw[decorate,decoration={brace,amplitude=4pt},thick,black!60]
    (-19,2) -- (-19,62);
  \foreach \j/\name in {0/x_0,1/x_1,2/x_2,3/x_3} {
    \pgfmathsetmacro{\x}{3+4*\j}
    \draw[wire] (\x,-8) -- (\x,{\x*10/18});
    \node[font=\scriptsize,anchor=north] at (\x,-8) {$\name$};
  }
  \node[note,anchor=north] at (9,-13) {inputs of the LUT\\(control bits of the multiplexer)};
  \draw[wire] (18,32) -- (32,32) node[right,font=\small] {$f(x_3,x_2,x_1,x_0) = t_{8x_3+4x_2+2x_1+x_0}$};
\end{tikzpicture}
\caption{A lookup table with four inputs: a 16-to-1 multiplexer whose data
inputs are constants. The inputs of the LUT select one of them.}
\label{fig:lut}
\end{figure}

What does such a LUT compute? The four input bits, read as a binary number
$i = 8x_3 + 4x_2 + 2x_1 + x_0 \in \{0, \dots, 15\}$, select the configuration
bit $t_i$. So the LUT computes the Boolean function
\[
  f\colon \{0,1\}^4 \to \{0,1\}, \qquad
  f(x_3, x_2, x_1, x_0) = t_{8x_3+4x_2+2x_1+x_0},
\]
and the 16 configuration bits are nothing but the last column of its truth
table. Every Boolean function of four variables has a truth table with 16 rows,
so a LUT can compute \emph{every} such function: there are $2^{16} = 65536$ of
them, one for each way to choose the configuration bits.

This is the whole trick of an FPGA. There are no AND or OR gates on the chip
waiting to be connected. A circuit that you draw with gates is translated into
Boolean functions of at most four inputs each, and every one of these functions
becomes a LUT with the right 16 bits. A function of more than four inputs is
split into several LUTs, the output of one feeding an input of the next.

\begin{exercise}
Which configuration bits $t_0, \dots, t_{15}$ make a LUT compute the
function $f(x_3,x_2,x_1,x_0) = x_3 \wedge x_2$? Which ones make it compute
$x_0 \xor x_1 \xor x_2 \xor x_3$?
\end{exercise}

\section{Flip-Flops, and Why the Clock Matters}
\label{sec:ff-intro}

A circuit of LUTs alone is combinational: its outputs depend only on its
current inputs. To store something we need flip-flops, and every logic cell has
one right behind its LUT. The output of the LUT can go out of the logic cell
directly, or into the flip-flop, which takes it over at the rising edge of a
clock.

On our board a crystal oscillator delivers a clock of 12\,MHz, that is
$12 \cdot 10^6$ rising edges per second, one every $83.3$\,ns. All flip-flops
in our designs use this one clock (or are meant to, see
Section~\ref{sec:step02-pins}). Between two edges, the signals have to travel
from the flip-flops through LUTs and wires to the next flip-flops and settle
there. This takes time, and so the tools tell us after placing the design how
fast the clock could be. For our larger designs this will be a real
limitation.

\section{The Board}
\label{sec:board}

The iCEBreaker (Figure~\ref{fig:board}) carries the FPGA, the 12\,MHz
oscillator, a flash memory that stores the configuration, and a USB chip. The
USB chip has two jobs: it writes the configuration into the flash, and it
provides a serial line between your computer and the FPGA, which we use from
Chapter~\ref{ch:step06} on.

\begin{figure}[ht]
\centering
\begin{tikzpicture}[node distance=6mm and 22mm]
  \node[chipbox,minimum width=34mm,minimum height=24mm] (fpga) {iCE40UP5K\\[-2pt]\footnotesize FPGA};
  \node[box,fill=black!5,left=of fpga,minimum width=24mm,minimum height=10mm,yshift=8mm] (usb) {USB chip\\[-2pt]\footnotesize FTDI FT2232H};
  \node[box,fill=black!5,left=of fpga,minimum width=24mm,minimum height=10mm,yshift=-8mm] (flash) {flash\\[-2pt]\footnotesize 16\,MiB};
  \node[box,fill=black!5,above=of fpga,minimum width=24mm] (osc) {oscillator\\[-2pt]\footnotesize 12\,MHz};
  \node[box,fill=yellow!12,right=12mm of fpga,minimum width=28mm,yshift=10mm] (pmod1) {PMOD 1A\\[-2pt]\footnotesize 7-segment display};
  \node[box,fill=yellow!12,right=12mm of fpga,minimum width=28mm,yshift=-10mm] (pmod2) {PMOD 2\\[-2pt]\footnotesize 3 buttons, 5 LEDs};
  \node[left=10mm of usb,note] (pc) {your\\computer};
  \draw[biarrow] (pc) -- node[above,note]{USB} (usb);
  \draw[arrow] (usb) -- node[left,note]{writes} (flash);
  \draw[arrow] (flash) -- node[above,note]{configuration} (fpga.west |- flash);
  \draw[biarrow] (usb) -- node[above,note]{serial} (fpga.west |- usb);
  \draw[arrow] (osc) -- node[right,note]{clock} (fpga);
  \draw[biarrow] (fpga.east |- pmod1) -- (pmod1);
  \draw[biarrow] (fpga.east |- pmod2) -- (pmod2);
\end{tikzpicture}
\caption{What is on the board, and how the parts are connected.}
\label{fig:board}
\end{figure}

The FPGA forgets its configuration when the power goes off. At power-up it
reads the configuration from the flash, and so the flash is what we program.

On the right edge of the board are connectors in the PMOD format, a simple
standard for small extension modules. On one we plug the module with three
buttons and five LEDs, on another the 7-segment display with two digits.

Which wire of the FPGA leads to which button is not something the FPGA knows.
The file \cc{icebreaker.pcf} at the top of the repository gives the pins of
the chip names that we can use in our designs. A few lines of it:
\begin{terminal}
# 12 MHz clock
set_io -nowarn CLK        35
...
# LEDs and Buttons (PMOD 2)
set_io -nowarn LED1       26
set_io -nowarn LED2       27
...
set_io -nowarn BTN1       20
set_io -nowarn BTN2       19
\end{terminal}
So pin 35 of the chip is called \cc{CLK}, and pin 20 \cc{BTN1}. When a design
has an input called \cc{BTN1}, it is connected to the first button.

\section{From a Description to a Configured Chip}
\label{sec:toolflow}

Between the SystemVerilog file and the configured chip there are four tools,
and it pays to know what each of them does. Each one reads a file and writes
the next (Figure~\ref{fig:toolflow}).

\begin{figure}[ht]
\centering
\begin{tikzpicture}[node distance=5mm and 7mm]
  \node[filebox] (sv) {\cc{test.sv}};
  \node[toolbox,right=of sv] (sv2v) {\textbf{sv2v}\\[-2pt]\footnotesize translate};
  \node[filebox,right=of sv2v] (v) {\cc{test.v}};
  \node[toolbox,right=of v] (yosys) {\textbf{yosys}\\[-2pt]\footnotesize synthesize};
  \node[filebox,right=of yosys] (json) {\cc{test.json}};
  \node[toolbox,below=9mm of json] (pnr) {\textbf{nextpnr}\\[-2pt]\footnotesize place \& route};
  \node[filebox,left=of pnr] (asc) {\cc{test.asc}};
  \node[toolbox,left=of asc] (pack) {\textbf{icepack}\\[-2pt]\footnotesize pack bits};
  \node[filebox,left=of pack] (bin) {\cc{test.bin}};
  \node[toolbox,left=of bin] (prog) {\textbf{iceprog}\\[-2pt]\footnotesize write flash};
  \node[filebox,right=of pnr] (pcf) {\cc{icebreaker.pcf}};
  \draw[arrow] (sv) -- (sv2v); \draw[arrow] (sv2v) -- (v);
  \draw[arrow] (v) -- (yosys); \draw[arrow] (yosys) -- (json);
  \draw[arrow] (json) -- (pnr); \draw[arrow] (pcf) -- (pnr);
  \draw[arrow] (pnr) -- (asc); \draw[arrow] (asc) -- (pack);
  \draw[arrow] (pack) -- (bin); \draw[arrow] (bin) -- (prog);
\end{tikzpicture}
\caption{The tools between a description in SystemVerilog and the
configuration in the flash of the board. \cc{make synth} runs everything up to
\cc{test.bin}, \cc{make prog} also \cc{iceprog}.}
\label{fig:toolflow}
\end{figure}

\begin{description}
\item[sv2v] translates SystemVerilog into plain Verilog. Verilog is the older
  language from which SystemVerilog grew; yosys understands only part of
  SystemVerilog, but all of Verilog. For our first designs the translation
  hardly changes anything; later, SystemVerilog features such as interfaces
  and packages are translated into Verilog that means the same.
\item[yosys] \term{synthesizes} the design: it turns the description into a
  \term{netlist}, a list of the building blocks of our FPGA (LUTs with their 16
  bits, flip-flops, memory blocks) and of the wires between them. It also
  simplifies the logic as far as it can. Nothing is placed on the chip yet.
\item[nextpnr] \term{places} every building block of the netlist on one of the
  physical blocks of the chip and \term{routes} the wires, that is, it chooses
  the switches of the routing so that the right blocks are connected. Here the
  pin names from \cc{icebreaker.pcf} come in. Then it computes how long the
  signals take, and how fast the clock may be.
\item[icepack] converts the result into the bit pattern that the FPGA loads at
  power-up, the \term{bitstream}.
\item[iceprog] writes the bitstream through USB into the flash and restarts
  the FPGA, which then loads it.
\end{description}

We do not call these tools by hand. Every step directory contains the same
\cc{Makefile} (it is a link to the one in \cc{01_logic_gates}), and its
targets do it for us. The design is always in \cc{test.sv}, with the top
module called \cc{test}.

Next to this path to the chip there is a second one that does not need the
board at all: \term{simulation}. The tool Verilator translates the design into
a C++ class, and a small C++ program, the \term{testbench}, sets the inputs of
this class step by step and records how all signals change. We look at these
recordings in a waveform viewer. Chapter~\ref{ch:step01} shows both paths for
our first design. Appendix~\ref{app:tools} lists all \cc{make} targets.

\solutions

\solution
For $x_3 \wedge x_2$, the output is 1 exactly when $8x_3 + 4x_2 + 2x_1 + x_0
\ge 12$, that is, for the indices 12, 13, 14, 15. So $t_{12} = t_{13} = t_{14}
= t_{15} = 1$ and all other bits are 0. Written from $t_{15}$ down to $t_0$, as
yosys does it (Section~\ref{sec:step01-yosys}):
\cc{1111000000000000}.

For $x_0 \xor x_1 \xor x_2 \xor x_3$, the output is 1 exactly when the
index $i$ has an odd number of ones in binary: $i = 1, 2, 4, 7, 8, 11, 13, 14$.
From $t_{15}$ down to $t_0$: \cc{0110100110010110}.
